\documentclass[sigconf,nonacm]{acmart}
\usepackage{tabularx}
\emergencystretch=1em

% Undergraduate CHI-style assignment template adapted from main.md.
% Replace every bracketed placeholder before submission.
\setcopyright{none}
\settopmatter{printacmref=false,printfolios=true}
\renewcommand\footnotetextcopyrightpermission[1]{}

\newcommand{\placeholderfigure}[2]{%
  \begin{figure}[tb]
    \centering
    \fbox{\parbox[c][1.15in][c]{0.88\columnwidth}{\centering\small #1\\[0.5em]
    \textit{Replace with a legible, original study artifact.}}}
    \caption{#2}
    \Description{Placeholder for an original, de-identified study artifact.}
  \end{figure}}

\newcommand{\guidance}[1]{\paragraph{Writing guidance.} \textit{#1}}
\newcommand{\instructornote}[1]{\paragraph{Instructor note.} \textit{#1}}
\newcommand{\commonmistakes}[1]{\paragraph{Common mistakes.} \textit{#1}}
\newcommand{\expectedlength}[1]{\paragraph{Expected length.} \textit{#1}}

\begin{document}

\title{Discovering, Establishing, and Designing for Requirements: An HCI Case Study}
\subtitle{Undergraduate CHI-Style Assignment Template}

\author{Anonymous Author(s)}
\affiliation{%
  \institution{Anonymous Institution}
  \city{Dhaka}
  \country{Bangladesh}}

\begin{abstract}
This report presents a complete requirements-to-conceptual-design account for
[system or service]. It documents how the team scoped a design problem,
gathered and analyzed evidence, established defensible requirements, and
translated those requirements into a conceptual design. Replace bracketed text
with a coherent account of one project; do not leave the abstract as a list of
activities. The finished abstract should state the target users, setting,
methods, principal findings, requirements, and design contribution in
150--200 words.
\end{abstract}

\keywords{human-centered design, requirements, data gathering, data analysis,
conceptual design, HCI education}

\maketitle

\section{Introduction}

\subsection{Problem, Users, and Setting}

[Name the setting] currently requires [primary users] to [goal], but this is
difficult because [observed or evidenced breakdown]. This project investigates
how an interactive system could support [bounded outcome] without assuming
that a new application is automatically the right answer. The unit of analysis
is [individual / team / service encounter], and the study concerns [location,
frequency, constraints, and stakeholders]. State what is known from preliminary
observation, prior research, or institutional context; clearly label conjecture
as conjecture.

The project is grounded in a human-centered process: discover the problem
space, gather evidence about people and practice, interpret that evidence,
negotiate requirements, and develop a conceptual design
\cite{preece2023interaction}. This sequence is iterative, not a claim that
requirements become permanently fixed after one study.

\begin{table}[tb]
  \caption{Project scope. Replace all bracketed entries.}
  \label{tab:scope}
  \begin{tabularx}{\columnwidth}{@{}lX@{}}
    \toprule
    Element & Project decision \\
    \midrule
    Primary users & [Who directly uses or is affected by the system?] \\
    Stakeholders & [Who administers, supports, regulates, or is indirectly affected?] \\
    Activity & [Observable task or practice under study] \\
    Setting & [Physical, social, organizational, and technical context] \\
    Boundary & [What is intentionally outside this assignment?] \\
    \bottomrule
  \end{tabularx}
\end{table}

\placeholderfigure{A context photograph, service-blueprint excerpt, or
stakeholder map for [setting]. Remove identifying information.}{[Replace this
caption.] Describe what the figure shows and why it matters.}

\guidance{Write in the past tense for completed study work and present tense
for stable context. Give one concrete example of the existing workflow.
Replace every bracketed placeholder, and cite sources for contextual claims.}

\instructornote{Assess scope discipline: a strong introduction identifies a
feasible population and activity, distinguishes stakeholders, and makes a
testable design opportunity visible.}

\subsection{Research Questions and Contribution}

The study asks: \textbf{RQ1:} How do [users] currently accomplish [goal]?
\textbf{RQ2:} Which needs, constraints, and opportunities should shape a future
design? \textbf{RQ3:} How can a conceptual design make the proposed interaction
understandable and accountable? The contribution is a documented,
evidence-traceable requirements set and concept for [system], not merely a
polished interface.

\expectedlength{1.5--2 pages, including the context figure and Table~\ref{tab:scope}.}
\commonmistakes{Opening with a solution; treating all people as ``the user'';
claiming needs without evidence; using research questions that can be answered
only by opinions about the proposed interface.}

\section{Discovering Requirements}

\subsection{Stakeholders, Assumptions, and Problem Space}

Requirements discovery identifies the people, activities, values, and
constraints that define a design space \cite{preece2023interaction}. We began
with a stakeholder inventory. Primary participants were [group]; secondary
stakeholders included [group]; and gatekeepers or domain experts included
[group]. Their interests were not assumed to align: for example, [participant
group] values [value], whereas [organization] must also satisfy [policy, safety,
cost, or privacy constraint].

\begin{table}[tb]
  \caption{Stakeholder analysis.}
  \label{tab:stakeholders}
  \begin{tabularx}{\columnwidth}{@{}l l X@{}}
    \toprule
    Stakeholder & Relationship & Interest, influence, and concern \\
    \midrule
    Primary user & Direct use & [Goal; capability; risk if the system fails] \\
    Secondary user & Indirect use & [Coordination need or downstream effect] \\
    Gatekeeper & Access/control & [Permission, policy, resource, or ethical concern] \\
    Domain expert & Knowledge source & [Practice knowledge and potential bias] \\
    \bottomrule
  \end{tabularx}
\end{table}

\guidance{Use short, concrete vignettes: ``When \ldots, [role] \ldots because
\ldots''. Mark statements as observation, participant account, document
evidence, or team inference. Show why the problem is consequential for people,
not just technically interesting.}

\instructornote{Look for a credible problem-space account. The best work
exposes tensions and uncertainty, then uses research to resolve them rather
than defending an early solution.}

\commonmistakes{Listing demographics instead of stakeholders' roles; treating
a persona as data; saying ``users need an app''; confusing a desired feature
with a human need; omitting organizational and accessibility constraints.}

\subsection{Current Practice and Opportunity Areas}

Describe a representative journey from trigger to outcome. For [persona role],
the activity begins when [trigger]. They then [step], coordinate with [person or
tool], and finish when [criterion]. Friction appears at [breakdown], where
[consequence]. Do not convert every complaint into a feature request:
distinguish a cause, an observed workaround, and a possible design opportunity.

Three preliminary opportunity areas were [opportunity A], [opportunity B], and
[opportunity C]. We carried these forward as questions to investigate, rather
than commitments to implement. Contextual inquiry is useful when work is
situated and tacit because it lets a researcher ask about decisions close to
the activity \cite{beyer1998contextual}.

\placeholderfigure{A stakeholder map or current-state journey. Use arrows to
show relationships and label tensions, handoffs, and pain points.}{[Replace this
caption.] Describe what the current practice artifact shows and why it matters.}

\expectedlength{2--2.5 pages. Include a stakeholder table and one
current-practice artifact.}

\section{Data Gathering}

\subsection{Method Rationale and Study Design}

We selected [interviews / observations / survey / diary study / document
analysis] because it can reveal [kind of evidence] in [context]. The method
complements rather than replaces [second method]: [method A] captures
[strength], while [method B] checks or broadens [strength]. Data-gathering
choices should be driven by questions, practical constraints, and
participants' wellbeing \cite{preece2023interaction}.

Participants were recruited through [channel] using [inclusion criteria]. We
sought variation in [experience, role, access need, or context], not statistical
representativeness. State the final sample honestly: $n=[\ ]$, with
[distribution]. Explain any attrition or access limitation and how it affects
the claims that follow.

We recorded initial assumptions in a risk register: [assumption 1], [assumption
2], and [assumption 3]. Each assumption was marked as either evidence-backed,
uncertain, or contradicted after fieldwork. This protects the project from
designing around the team's own habits rather than participants' practices.

\subsection{Procedure, Ethics, and Data Management}

Each session lasted approximately [duration]. After consent, participants
[procedure steps]. We gathered [audio, notes, artifacts, survey responses] and
stored them in [approved location] with [pseudonyms/access controls].
Participation was voluntary; participants could skip questions or withdraw
[state applicable policy]. Avoid collecting identifiers that the analysis does
not require.

The protocol included open prompts such as ``Tell me about the last time you
[activity]'' and follow-ups such as ``What happened next?'' Questions about
actual episodes reduce speculative answers. Pilot [number] session(s) revealed
[change], so we revised [prompt/order/material].

\begin{table*}[tb]
  \caption{Data-gathering protocol and evidence produced.}
  \label{tab:protocol}
  \begin{tabularx}{\textwidth}{@{}l X X@{}}
    \toprule
    Stage & Researcher action & Evidence and safeguard \\
    \midrule
    Recruitment & [Invite and screen] & [Eligibility; voluntary participation] \\
    Opening & [Consent and context] & [Information sheet; permission to record] \\
    Core activity & [Observe / interview / task] & [Notes, quotes, artifacts] \\
    Closing & [Debrief and compensate] & [Withdrawal/contact information] \\
    After session & [De-identify and store] & [Pseudonym; access control; retention] \\
    \bottomrule
  \end{tabularx}
\end{table*}

\placeholderfigure{A data-gathering timeline: recruitment, consent, session
activities, debrief, and secure handling of data.}{[Replace this caption.]
Describe the data-gathering timeline and why it matters.}

\guidance{Write the procedure chronologically. Separate what you planned from
what happened. Include one or two actual question stems, but place the full
guide in Appendix~\ref{app:materials}.}

\instructornote{Grade the fit between research questions and method. Ethical
care, sampling transparency, and a plausible procedure matter more than an
unnecessarily large sample.}

\expectedlength{2.5--3 pages. Include enough protocol detail for a peer to
reproduce the study responsibly.}
\commonmistakes{Claiming a convenience sample represents everyone; asking
leading questions; describing consent as a signature only; failing to state
what data were recorded; reporting a survey without its questions, scale,
recruitment, or response count.}

\section{Data Analysis, Interpretation, and Presentation}

\subsection{Preparation and Analytic Approach}

We analyzed [data types] using [thematic analysis / content analysis /
descriptive statistics / affinity diagramming]. First, [transcription, cleaning,
or familiarization step] prepared the corpus. We then assigned concise codes to
meaningful segments, compared codes across participants and contexts, and
grouped related codes into candidate themes. Themes were reviewed against raw
evidence and renamed to express an interpretive pattern, not a topic label
\cite{braun2006thematic}.

For transparency, each finding below has an evidence trail: source material
$\rightarrow$ code $\rightarrow$ theme $\rightarrow$ design implication.
[Number] team member(s) coded the material. Explain how disagreements were
discussed, how a codebook changed, and how negative cases were treated. Do not
report inter-rater reliability unless your analysis design actually used it and
you can explain it.

\subsection{Findings}

\paragraph{Theme 1: [interpretive statement].} Participants described
[pattern]. For example, P[ID] said, ``[short, anonymized quotation].'' This
matters because [interpretation tied to context], not merely because the quote
is vivid.

\paragraph{Theme 2: [interpretive statement].} Evidence from
[participant/artifact/observation] showed [pattern], including the
counterexample [case]. This suggests [implication] only within [boundary].

\paragraph{Theme 3: [interpretive statement].} [Explain pattern, variation,
and consequence.] Quantitative results, if used, should give denominator,
measure, and uncertainty where appropriate: ``[x]/[n] respondents selected
[response]'' is more informative than ``most users agreed.''

\begin{table*}[tb]
  \caption{Evidence-to-interpretation traceability matrix. Use anonymized identifiers.}
  \label{tab:evidence}
  \begin{tabularx}{\textwidth}{@{}X X X X@{}}
    \toprule
    Evidence source & Code / observation & Theme & Interpretation and bounded implication \\
    \midrule
    P01 interview, [time] & [Concise segment meaning] & [Interpretive theme] & [Why this pattern matters; do not write a feature yet] \\
    Observation [site], [time] & [Action / workaround / breakdown] & [Theme] & [Contextual explanation and exception] \\
    Artifact [ID] & [Documented constraint] & [Theme] & [Constraint on future requirement] \\
    \bottomrule
  \end{tabularx}
\end{table*}

\placeholderfigure{An affinity-map excerpt, coded journey, or chart. Show
enough labels to make the transformation from data to theme understandable.}{[Replace
this caption.] Describe the analysis artifact and why it matters.}

\guidance{Use descriptive headings that make an argument (for example,
``Coordination breaks down at handoffs''), not headings such as ``Interview
question 3.'' Present evidence first, interpretation second, and implication
third. Retain a de-identified audit trail in Appendix~\ref{app:audit}.}

\instructornote{Strong analysis is traceable and selective. It explains why
patterns matter to participants' activity, acknowledges variation, and avoids
claiming universality from a small qualitative sample.}

\expectedlength{3--3.5 pages. This is often the intellectual center of the report.}
\commonmistakes{Presenting quotes as self-explanatory; turning every code into a
theme; erasing contradictory evidence; reporting percentages with a tiny or
unstated denominator; mixing a result and proposed feature in the same sentence.}

\section{Establishing Requirements}

\subsection{From Findings to Requirements}

Requirements state what a system and its surrounding service must enable or
respect; they should not prematurely prescribe a screen layout
\cite{preece2023interaction}. We translated each supported finding into a
requirement by naming the actor, condition, action or quality, and success
criterion. Requirements were checked with [participants, stakeholders, expert,
or literature] through [walkthrough, prioritization activity, or review].

Functional requirements describe supported actions. Data requirements identify
the information needed, its source, ownership, sensitivity, and lifecycle.
Non-functional requirements express qualities such as accessibility, privacy,
reliability, learnability, and performance. Organizational requirements
describe policy, staffing, integration, or governance conditions. The four
categories can interact: an apparently simple reminder feature may create data,
consent, and notification constraints.

\placeholderfigure{A requirements traceability diagram linking findings,
requirement IDs, and later conceptual-design elements.}{[Replace this caption.]
Describe the traceability diagram and why it matters.}

\subsection{Prioritization and Acceptance Criteria}

We prioritized requirements using [MoSCoW / value-risk matrix / stakeholder
workshop]. A ``Must'' item is essential for the defined release or study
scenario, not merely popular. Priority reflects evidence, consequence of
absence, implementation dependencies, and equity impacts. For each
high-priority requirement, specify a testable acceptance criterion: ``Given
[condition], [user] can [action] with [measure] without [harm].''

Requirements remain provisional when evidence is thin. R[ID] is flagged for
validation because [uncertainty]. This honest uncertainty is a useful design
result: it directs the next study rather than disguising a decision as fact.

\begin{table*}[tb]
  \caption{Requirements catalogue. Add rows; retain stable IDs for traceability.}
  \label{tab:requirements}
  \begin{tabularx}{\textwidth}{@{}l l X X l@{}}
    \toprule
    ID & Type & Requirement & Evidence / rationale & Priority \\
    \midrule
    R1 & Functional & The system shall enable [role] to [action] when [condition]. & Theme T[ID]; P[ID] & Must \\
    R2 & Data & The system shall collect/display [data] only for [purpose] and [retention rule]. & Constraint [ID] & Must \\
    R3 & Quality & A [user with access need] shall complete [task] using [support/criterion]. & Observation [ID] & Should \\
    R4 & Organizational & [Responsible role] shall be able to [govern/support action]. & Stakeholder review & Could \\
    \bottomrule
  \end{tabularx}
\end{table*}

\guidance{Use active, testable wording: ``The system shall enable [role] to
\ldots'' or ``[Role] needs to \ldots when \ldots''. Include the rationale and
evidence ID beside every important requirement. Do not invent requirements
merely to fill categories.}

\instructornote{Assess requirement quality, not quantity. A good table supports
later design decisions, identifies trade-offs, and makes it possible to
challenge an unsupported item.}

\expectedlength{2.5--3 pages, including a requirements table with at least
8--12 substantive entries.}
\commonmistakes{Writing features instead of needs; using vague terms such as
``easy'' without a criterion; prioritizing by personal preference; omitting
privacy and accessibility; presenting all requirements as equally important;
failing to trace requirements to evidence.}

\section{Conceptual Design}

\subsection{Design Vision and Conceptual Model}

The conceptual design proposes [system name], a [type of interactive
product/service] that helps [role] [goal] in [context]. Its central design
principle is [principle], derived from R[ID] and T[ID]. The system model is
[object/action model]: users work with [key objects] and can [key actions].
Define these terms in language participants recognize; a conceptual model
should help users predict what actions are possible and what outcomes follow
\cite{norman2013design}.

The proposed interaction is [direct manipulation / conversation / form fill-in
/ notification / embodied interaction / hybrid]. Explain why this style fits
the activity, environment, devices, attention demands, and access needs. State
deliberate exclusions, such as ``The system does not make [high-stakes
decision]; it supports [human role] with [information or coordination].''

\placeholderfigure{A conceptual-model diagram or task flow. Use named objects,
actions, feedback, and exception paths; annotate with R1--R[n].}{[Replace this
caption.] Describe the conceptual model and why it matters.}

\subsection{Scenario and Task Flow}

\paragraph{Scenario: [name].} [Persona/role] begins when [trigger]. They see
[information] and choose [action]. The system then [response], making [status,
consequence, or uncertainty] visible. If [exception], the system [recovery or
escalation]. The scenario ends when [success criterion]. Annotate each step with
the requirement(s) it serves.

\subsection{Design Alternatives and Trade-Offs}

We considered [alternative A] and [alternative B]. Alternative A was rejected
or deferred because [evidence-based trade-off]; alternative B was retained
because [reason]. Discuss effects on workload, privacy, accessibility, control,
error recovery, and organizational practice. Participatory design is especially
valuable when those affected can shape consequential design choices
\cite{muller1993participatory}.

\begin{table*}[tb]
  \caption{Requirement-to-conceptual-design mapping.}
  \label{tab:designmapping}
  \begin{tabularx}{\textwidth}{@{}l X X@{}}
    \toprule
    Requirement & Conceptual-design response & How it will be checked \\
    \midrule
    R1 & [Object, action, and feedback] & [Scenario / usability task] \\
    R2 & [Data boundary or consent mechanism] & [Policy review / walkthrough] \\
    R3 & [Accessibility or resilience support] & [Inclusive test condition] \\
    R4 & [Administration or service process] & [Stakeholder review] \\
    \bottomrule
  \end{tabularx}
\end{table*}

\guidance{Describe the interaction before drawing it. Name objects consistently
across scenario, diagram, and prototype. A conceptual design may use
low-fidelity sketches, but each sketch must answer a requirement or reveal a
decision.}

\instructornote{Look for a design that is intelligible before visual polish.
Strong submissions show traceability, explicit trade-offs, and appropriate
limits on what the technology decides.}

\expectedlength{3--3.5 pages. Include a model/flow and a mapping from
requirements to design.}
\commonmistakes{Jumping directly to annotated screens; calling a feature list a
conceptual model; ignoring exceptions; treating automation as neutral; using a
persona without linking it to study evidence; failing to explain why an
interaction style fits the context.}

\section{Evaluation Plan}

\subsection{Evaluation Questions and Methods}

The next iteration will evaluate whether the conceptual design supports the
priority requirements, rather than whether participants simply like it.
Evaluation questions are: (1) Can [role] understand the model and complete
[task]? (2) Does [feature/flow] create confusion, workload, exclusion, or
privacy concern? (3) Does the design fit [organizational/workflow constraint]?
A moderated think-aloud walkthrough with [target participants] will test
comprehension and task flow; a stakeholder review will test operational
feasibility.

The prototype will be [paper / clickable low-fidelity / service enactment],
intentionally matched to the questions. A paper prototype can test sequence,
labels, and feedback but cannot establish real-world reliability. State what
the prototype cannot validly evaluate.

\begin{table*}[tb]
  \caption{Evaluation plan tied to requirements.}
  \label{tab:evaluation}
  \begin{tabularx}{\textwidth}{@{}X X X@{}}
    \toprule
    Task / requirement ID & Method and evidence & Decision rule \\
    \midrule
    {[Task], R1} & [Think-aloud; completion, errors, quotation] & [Revise task if \ldots] \\
    {[Task], R2} & [Walkthrough; data-flow understanding] & [Escalate if \ldots] \\
    {[Task], R3} & [Inclusive session; barrier observed] & [Revise if \ldots] \\
    \bottomrule
  \end{tabularx}
\end{table*}

\placeholderfigure{An evaluation journey showing participant tasks,
observation points, measures, and revision decisions.}{[Replace this caption.]
Describe the evaluation journey and why it matters.}

\subsection{Success, Risk, and Iteration}

Success criteria derive from Table~\ref{tab:requirements}. For example,
[x]/[n] participants should [task outcome], and no participant should encounter
an unaddressed [severity] breakdown related to R[ID]. Qualitative observations
will be coded for misunderstandings, workarounds, breakdowns, and unanticipated
value. Findings will lead to [revision rule], while high-risk concerns such as
[privacy/safety/equity concern] pause implementation until reviewed.

\guidance{For each task, specify the starting condition, prompt, expected
outcome, evidence collected, and decision triggered by a failure. Avoid a
generic ``we will conduct usability testing'' sentence.}

\instructornote{Reward a realistic, requirement-driven evaluation plan. Its
purpose is to reduce uncertainty and guide revision, not to certify a design as
finished.}

\expectedlength{1.5--2 pages. This is a planned study: use future tense and
avoid inventing results.}
\commonmistakes{Using satisfaction as the only outcome; testing only happy
paths; claiming a prototype tests what it cannot; writing success measures
unrelated to requirements; overlooking participant risk or accessibility in
the evaluation.}

\section{Discussion and Conclusion}

\subsection{What the Process Established}

This project established that [main supported finding] shapes the design of
[system] because [reason]. The resulting requirements prioritize [two or three
priorities], and the conceptual design responds through [named design moves].
The value of the work lies in the chain from situated evidence to a design that
can be evaluated, rather than in claiming that one small study has discovered
universal user needs.

\begin{table}[tb]
  \caption{Bounded claims and next steps.}
  \label{tab:claims}
  \begin{tabularx}{\columnwidth}{@{}X X X@{}}
    \toprule
    Supported claim & Boundary / limitation & Next evidence or decision \\
    \midrule
    {[Claim grounded in T/R IDs]} & [Population, setting, or method boundary] & [Validation or design action] \\
    {[Claim grounded in T/R IDs]} & [Uncertainty or negative case] & [Stakeholder decision] \\
    \bottomrule
  \end{tabularx}
\end{table}

\subsection{Limitations and Next Steps}

The study is limited by [sample, setting, duration, access, method, or
researcher-position limitation]. Consequently, findings should be interpreted
as [bounded claim]. The next iteration should [specific evidence-gathering or
evaluation action], particularly to examine [uncertain requirement or
underserved stakeholder]. Reflect on how the team's position, access, and
assumptions may have influenced what participants shared and how evidence was
interpreted.

\placeholderfigure{A concise roadmap from the completed study to the next
validation, prototype iteration, and stakeholder decision.}{[Replace this
caption.] Describe the roadmap and why it matters.}

\subsection{Conclusion}

Conclude in one paragraph by restating the problem, evidence-based contribution,
and immediate next design decision. Do not introduce new claims here. The final
sentence should make clear what a future designer or researcher can do
differently because of this report.

\guidance{Use this section to make a bounded argument. Name both the
contribution and its limits. A candid limitation increases credibility when it
is connected to a concrete next step.}

\instructornote{The conclusion should demonstrate reflective HCI practice:
evidence, consequence, uncertainty, and a thoughtful path for iteration.}

\expectedlength{1--1.5 pages.}
\commonmistakes{Repeating the abstract word-for-word; claiming proof from a
small study; hiding limitations; proposing a new feature in the conclusion;
ending with a generic statement that technology will improve life.}

\begin{thebibliography}{5}

\bibitem{beyer1998contextual}
Hugh Beyer and Karen Holtzblatt. 1998.
\newblock \emph{Contextual Design: Defining Customer-Centered Systems}.
\newblock Morgan Kaufmann Publishers, San Francisco, CA.

\bibitem{braun2006thematic}
Virginia Braun and Victoria Clarke. 2006.
\newblock Using thematic analysis in psychology.
\newblock \emph{Qualitative Research in Psychology} 3, 2 (2006), 77--101.

\bibitem{muller1993participatory}
Michael J. Muller. 1993.
\newblock Participatory design: The third space in HCI.
\newblock In \emph{The Human--Computer Interaction Handbook}. 1051--1068.

\bibitem{norman2013design}
Donald A. Norman. 2013.
\newblock \emph{The Design of Everyday Things}, revised and expanded ed.
\newblock MIT Press, Cambridge, MA.

\bibitem{preece2023interaction}
Jennifer Preece, Yvonne Rogers, and Helen Sharp. 2023.
\newblock \emph{Interaction Design: Beyond Human--Computer Interaction}, 6th ed.
\newblock John Wiley \& Sons.

\end{thebibliography}

\appendix

\section{Study Materials}
\label{app:materials}

\subsection{Participant Information and Consent Checklist}

Before beginning, provide: the study purpose; what participation involves;
duration; recordings and data handling; foreseeable discomforts; benefits, if
any; voluntary participation and withdrawal process; contact details; and the
applicable approval statement. Do not include a participant's name, signature,
contact information, or other identifying data in this submitted report.

\subsection{Semi-Structured Interview Guide}

\begin{enumerate}
  \item Tell me about the last time you [performed the focal activity]. What
  happened from beginning to end?
  \item What parts were easy, difficult, slow, risky, or frustrating? Can you
  show me an example?
  \item Which people, tools, documents, and information did you use? What
  happened when one was unavailable?
  \item What workarounds have you developed, and what do they cost you or
  others?
  \item If a new system supported this activity, what would it need to respect
  about your practice, privacy, access, or control?
  \item Is there anything important we did not ask? Which of your answers should
  we treat cautiously?
\end{enumerate}

\paragraph{Interviewer reminders.} Ask neutral follow-ups (``Can you say
more?''); do not promise a feature; do not pressure participants to disclose
sensitive information; record analytic memos separately from factual notes.

\instructornote{Check that the guide could plausibly generate the evidence used
in the Findings section and that it treats participants with respect.}
\expectedlength{1--2 pages, only as needed. This appendix supports the method;
it does not replace the method description.}
\commonmistakes{Including identifiable data; presenting leading questions;
omitting a withdrawal/debrief procedure; pasting a raw transcript without
purpose or permission.}

\section{Analysis Audit Trail}
\label{app:audit}

\paragraph{Reflexive memo.} [Write 100--200 words on how the researcher's role,
expectations, relationship to the setting, language, or analytic decisions may
have shaped this account. State one action taken to examine that influence.]

\begin{table}[tb]
  \caption{Example de-identified analytic audit trail.}
  \label{tab:audit}
  \begin{tabularx}{\columnwidth}{@{}X X X@{}}
    \toprule
    Evidence / initial code & Analytic decision & Theme or requirement link \\
    \midrule
    {[De-identified excerpt or observation; concise code]} & [How interpretation changed] & [T/R ID] \\
    {[Negative or divergent case; concise code]} & [Why it was retained] & [Boundary / revision] \\
    \bottomrule
  \end{tabularx}
\end{table}

\instructornote{An audit trail makes reasoning inspectable. It is not a demand
for mechanical objectivity; it shows disciplined interpretation.}
\expectedlength{1--2 pages. Include representative evidence, not every note.}
\commonmistakes{Treating codes as objective facts; exposing identifiable
quotations; claiming triangulation merely because multiple data types exist;
omitting a negative case or analytic uncertainty.}

\end{document}
